\chapter*{Streszczenie}
Kompilator jest niezbędnym narzędziem przy tworzeniu programów komputerowych,
które efektywnie wykorzystują zasoby sprzętowe systemu. Bywa on jednak
traktowany jedynie jako tłumacz kodu, z~pominięciem jego możliwości
optymalizacyjnych, choć sposób tłumaczenia kodu źródłowego na instrukcje
maszynowe wpływa bezpośrednio na wydajność programu.

Zagadnienie to ma szczególne znaczenie w~systemach
wbudowanych, w~których poprawa wydajności oznacza często niższy
koszt i~mniejszy pobór energii. W~prezentowanej pracy zaprojektowano system
wbudowany oparty na platformie Raspberry Pi 4B z~systemem Linux, przeznaczony do
badania możliwości optymalizacyjnych kompilatorów GCC i~Clang. Jako algorytm
testowy wybrano szybką transformatę Fouriera, dla której zaimplementowano
w~języku C jedenaście wariantów, zarówno sekwencyjnych jaki~wielowątkowych. 

W~pracy zawarto również opis systemów operacyjnych w~zastosowaniach wbudowanych, podstaw
transformaty Fouriera oraz budowy kompilatora i~stosowanych przez niego
optymalizacji. Badania objęły jedenaście zestawów flag kompilacji i~trzy
rozmiary danych wejściowych, a~oprócz czasu wykonania analizowano liczniki
sprzętowe, rozmiar kodu wynikowego oraz dokładność obliczeń.

Największy zysk
przyniosło odejście od domyślnego poziomu \texttt{-O0}, które skróciło czas
wykonania od 2,18 do 5,11 raza, podczas gdy dalszy dobór flag dawał zwykle
około 12\%. Dla dużych zbiorów danych o~czasie wykonania decydował przede
wszystkim wzorzec dostępu do pamięci, dlatego wybór schematu obliczeń okazał
się równie istotny jak poziom optymalizacji. Ręczna wektoryzacja z~użyciem
instrukcji NEON oraz zrównoleglenie obliczeń przynosiły wyraźne korzyści tylko
w~określonych zakresach rozmiaru danych. Kompilator nie skorygował przy tym
błędów numerycznych wynikających z~konstrukcji algorytmu. Uzyskane wyniki
wskazują, że optymalizację programu należy zaczynać od doboru algorytmu
i~organizacji danych, a~ustawienia kompilacji traktować jako ich uzupełnienie.

\bigskip

\noindent\textbf{Słowa kluczowe:} optymalizacja kompilacji, kompilatory GCC i~Clang, szybka transformata Fouriera, systemy wbudowane, system operacyjny Linux

\bigskip

\noindent\textbf{Dziedzina nauki i techniki zgodna z OECD} Nauki inżynieryjne i techniczne, Elektrotechnika, elektronika i inżynieria informatyczna,
